\subsection{The tensor algebra and its two scalar embeddings}\label{sec:dpr:tensor}
The original Diamond--Posen ring-switching compiler
is formulated over $A=L\otimes_KL$, not merely
over a matrix transpose. We reconstruct the
specialization $K=\F_2$ of Construction~4.1
in~\cite{DiamondPosen2024}. Write
$n=2^\ell$, $L=\F_{2^n}$, and use the
binary basis $(\beta_v)_{v\in I_\ell}$ from Chapter~6.

\begin{definition}[Two embeddings]\label{def:dpr:embeddings}
Let $A:=L\otimes_{\F_2}L$ be the tensor product
with multiplication
$(a\otimes b)(c\otimes d)=ac\otimes bd$
and identity $1\otimes1$. Define
\[
 \varphi_0(a):=a\otimes1,\qquad
 \varphi_1(a):=1\otimes a.
\]
For $x=\sum_{i}a_i\otimes b_i$ define
the multiplication map
$\mu(x):=\sum_i a_ib_i\in L$.
\end{definition}

\begin{lemma}[The multiplication map is well defined]\label{lem:dpr:mu-defined}
The map $\mu:A\to L$ is an $\F_2$-linear map,
and $\mu(\varphi_0(a))=\mu(\varphi_1(a))=a$.
\end{lemma}
\begin{proof}
The map $(a,b)\mapsto ab$ is additive
in each argument and balanced over $\F_2$.
By the universal property of the tensor
product it induces a unique $\F_2$-linear
map $\mu$ with $\mu(a\otimes b)=ab$.
Evaluating at $a\otimes1$ and $1\otimes a$
gives the two identities.
\end{proof}

\begin{lemma}[Two commuting field embeddings]\label{lem:dpr:phi-embed}
Both $\varphi_0,\varphi_1:L\to A$
are injective unital ring homomorphisms,
and $\varphi_0(a)\varphi_1(b)=a\otimes b$.
\end{lemma}
\begin{proof}
Additivity, multiplicativity and preservation
of the unit follow from the defining tensor
relations and multiplication rule.
The formula for the product is
$(a\otimes1)(1\otimes b)=a\otimes b$.
If $\varphi_i(a)=0$, Lemma~\ref{lem:dpr:mu-defined}
implies $a=\mu(\varphi_i(a))=0$,
so both maps are injective.
The two images commute because $A$
is a commutative algebra.
\end{proof}

\begin{definition}[Column coordinates in the tensor algebra]\label{def:dpr:column}
For $x\in A$ define $\operatorname{Col}(x)_v\in L$
as the unique coefficients satisfying
\[
 \boxed{x=\sum_{v\in I_\ell}
   \varphi_0(\operatorname{Col}(x)_v)
   \varphi_1(\beta_v).}
\]
\end{definition}

\begin{lemma}[Existence and uniqueness of tensor columns]\label{lem:dpr:column-basis}
The family $(\varphi_1(\beta_v))_v$
is an $L$-basis of $A$ with scalar action
through $\varphi_0$, so
Definition~\ref{def:dpr:column} is well defined.
\end{lemma}
\begin{proof}
Every $b\in L$ has a unique expansion
$b=\sum_vd_v(b)\beta_v$ with
$d_v(b)\in\F_2$. For a pure tensor,
\[
 a\otimes b
 =\sum_vd_v(b)\,(a\otimes\beta_v)
 =\sum_v\varphi_0(ad_v(b))\varphi_1(\beta_v).
\]
Every tensor is a finite sum of pure tensors,
so the proposed family spans $A$ over the
left $L$-action.
For independence, let
$\sum_va_v\otimes\beta_v=0$.
For each $u$, apply the well-defined
linear map $\mathrm{id}_L\otimes d_u:A\to
L\otimes_{\F_2}\F_2\simeq L$;
the result is
$\sum_va_vd_u(\beta_v)=a_u=0$.
All coefficients vanish.
\end{proof}

\begin{lemma}[Columns of the two embeddings]\label{lem:dpr:embedding-columns}
For all $a,b\in L$,
\[
 \operatorname{Col}(\varphi_0(a))_v=a\,d_v(1),
 \qquad
 \operatorname{Col}(\varphi_1(b))_v=d_v(b).
\]
In particular
$\operatorname{Col}(\varphi_1(b))$
has entries in $\F_2$.
\end{lemma}
\begin{proof}
Use the basis expansions
$1=\sum_vd_v(1)\beta_v$ and
$b=\sum_vd_v(b)\beta_v$
in the second tensor factor,
then compare the unique
left-$L$ coefficients from
Lemma~\ref{lem:dpr:column-basis}.
\end{proof}

\begin{lemma}[Left scalar multiplication acts columnwise]\label{lem:dpr:column-scaling}
For all $a\in L$ and $x\in A$,
\[
 \operatorname{Col}(\varphi_0(a)x)_v
  =a\,\operatorname{Col}(x)_v.
\]
\end{lemma}
\begin{proof}
Insert the unique column expansion of
$x$, multiply each term by
$\varphi_0(a)$, use multiplicativity
of $\varphi_0$, and compare coefficients
in the left-$L$ basis.
\end{proof}

\begin{definition}[Weighted column contraction]\label{def:dpr:contraction}
For $q=(q_v)_v\in L^{I_\ell}$ define
\[
 C_q:A\to L,\qquad
 C_q(x):=\sum_vq_v\operatorname{Col}(x)_v.
\]
\end{definition}

\begin{lemma}[Contraction is linear in the left scalar action]\label{lem:dpr:contraction-linear}
For $a\in L$ and $x,y\in A$,
$C_q(\varphi_0(a)x+y)=aC_q(x)+C_q(y)$.
\end{lemma}
\begin{proof}
Every coordinate map $\operatorname{Col}$
is additive, and
Lemma~\ref{lem:dpr:column-scaling}
shows compatibility with the left
scalar action. Multiply by $q_v$
and sum over $v$.
\end{proof}

\begin{lemma}[Ordinary multiplication loses tensor information]\label{lem:dpr:multiplication-loss}
The map $\mu:A\to L$ is not injective
when $n>1$. More precisely, for any
$a\in L\setminus\F_2$ the tensor
$\varphi_0(a)+\varphi_1(a)$ is nonzero
but has image zero under $\mu$.
\end{lemma}
\begin{proof}
By Lemma~\ref{lem:dpr:mu-defined},
$\mu(\varphi_0(a)+\varphi_1(a))=a+a=0$.
Suppose the tensor were zero.
Its $v$-th column would satisfy
$a\,d_v(1)+d_v(a)=0$ for every $v$,
by Lemma~\ref{lem:dpr:embedding-columns}.
As $1\ne0$, some coordinate $d_v(1)=1$.
For that $v$ the equality gives
$a=d_v(a)\in\F_2$, a contradiction.
\end{proof}

\begin{remark}[Provenance and scope]
The two embeddings and the distinction
between columns and the multiplication map
are essential to Diamond--Posen's actual
ring-switching Construction~4.1.
The binary transposition map $\mathcal T$
of Chapter~7 is a related coordinate
operation, but it is not a substitute for
the entire tensor-algebra sumcheck.
\end{remark}
